\subsection{The Complex Representation Architecture (\texttt{complex})}

Python has a built-in \texttt{complex} type. A complex number contains two floating-point parts:

\begin{itemize}
    \item a real part;
    \item an imaginary part.
\end{itemize}

In Python source code, the imaginary component is written with \texttt{j}.

\begin{verbatim}
signal = 3 + 4j

print(f"signal: {signal}")
print(f"type:   {type(signal)}")

print()
print("selected complex methods/attributes:")
print(f"{'real':<15} {'real' in dir(signal)}")
print(f"{'imag':<15} {'imag' in dir(signal)}")
print(f"{'conjugate':<15} {'conjugate' in dir(signal)}")
print(f"{'__add__':<15} {'__add__' in dir(signal)}")
print(f"{'__mul__':<15} {'__mul__' in dir(signal)}")

print()
print(f"real part:      {signal.real}")
print(f"imaginary part: {signal.imag}")
print(f"conjugate:      {signal.conjugate()}")
\end{verbatim}

Possible output:

\begin{verbatim}
signal: (3+4j)
type:   <class 'complex'>

selected complex methods/attributes:
real            True
imag            True
conjugate       True
__add__         True
__mul__         True

real part:      3.0
imaginary part: 4.0
conjugate:      (3-4j)
\end{verbatim}

The value \texttt{3 + 4j} is one object of type \texttt{complex}. It is not a pair of separate Python variables. The real and imaginary components are accessible through attributes.

\subsubsection{Language-Native Integration via the Mathematical Imaginary Component \texttt{j} Literal}

Python uses \texttt{j} for the imaginary unit. This follows the convention used in engineering, where \texttt{i} is often already used for electric current.

The imaginary literal must have a numeric coefficient:

\begin{verbatim}
z1 = 4j
z2 = 3 + 4j
z3 = -2j
z4 = 0.5j

print(f"z1: {z1}")
print(f"z2: {z2}")
print(f"z3: {z3}")
print(f"z4: {z4}")

print()
print(f"type(z1): {type(z1)}")
print(f"type(z2): {type(z2)}")
\end{verbatim}

Possible output:

\begin{verbatim}
z1: 4j
z2: (3+4j)
z3: (-0-2j)
z4: 0.5j

type(z1): <class 'complex'>
type(z2): <class 'complex'>
\end{verbatim}

The text \texttt{4j} is a complex literal. The text \texttt{j} alone is not an imaginary literal. It is treated as an ordinary name.

\begin{verbatim}
try:
    print(j)
except NameError as err:
    print("name lookup failed")
    print(f"message: {err}")
\end{verbatim}

Possible output:

\begin{verbatim}
name lookup failed
message: name 'j' is not defined
\end{verbatim}

This works only after binding \texttt{j} as a normal variable:

\begin{verbatim}
j = "Jerry"

print(j)
print(type(j))
\end{verbatim}

Possible output:

\begin{verbatim}
Jerry
<class 'str'>
\end{verbatim}

So the rule is:

\begin{quote}
\texttt{4j} is an imaginary numeric literal. Plain \texttt{j} is just a name.
\end{quote}

Uppercase \texttt{J} also works:

\begin{verbatim}
z1 = 4j
z2 = 4J

print(f"z1: {z1}")
print(f"z2: {z2}")
print(f"z1 == z2: {z1 == z2}")
\end{verbatim}

Possible output:

\begin{verbatim}
z1: 4j
z2: 4j
z1 == z2: True
\end{verbatim}

Complex arithmetic is built into the language:

\begin{verbatim}
a = 3 + 4j
b = 2 - 1j

print(f"a: {a}")
print(f"b: {b}")

print()
print(f"a + b: {a + b}")
print(f"a - b: {a - b}")
print(f"a * b: {a * b}")
print(f"a / b: {a / b}")
\end{verbatim}

Possible output:

\begin{verbatim}
a: (3+4j)
b: (2-1j)

a + b: (5+3j)
a - b: (1+5j)
a * b: (10+5j)
a / b: (0.4+2.2j)
\end{verbatim}

The operators are dispatched to complex-number behavior. The name \texttt{a} is still only a binding to a complex object.

\subsubsection{Component Access Mechanisms: Extracting \texttt{.real} and \texttt{.imag} Floating-Point Parts}

A complex object exposes its two components through \texttt{.real} and \texttt{.imag}.

\begin{verbatim}
position = 12.5 - 7.25j

real_part = position.real
imag_part = position.imag

print(f"position:  {position}")
print(f"real_part: {real_part}")
print(f"imag_part: {imag_part}")

print()
print(f"type(position):  {type(position)}")
print(f"type(real_part): {type(real_part)}")
print(f"type(imag_part): {type(imag_part)}")
\end{verbatim}

Possible output:

\begin{verbatim}
position:  (12.5-7.25j)
real_part: 12.5
imag_part: -7.25

type(position):  <class 'complex'>
type(real_part): <class 'float'>
type(imag_part): <class 'float'>
\end{verbatim}

The components are floating-point values. Even when the source text uses whole numbers, the extracted parts are floats:

\begin{verbatim}
z = 3 + 4j

print(f"z.real: {z.real}")
print(f"z.imag: {z.imag}")

print()
print(f"type(z.real): {type(z.real)}")
print(f"type(z.imag): {type(z.imag)}")
\end{verbatim}

Possible output:

\begin{verbatim}
z.real: 3.0
z.imag: 4.0

type(z.real): <class 'float'>
type(z.imag): <class 'float'>
\end{verbatim}

A complex number with no explicit real part has real part \texttt{0.0}:

\begin{verbatim}
z = 42j

print(f"z:      {z}")
print(f"real:   {z.real}")
print(f"imag:   {z.imag}")
print(f"type:   {type(z)}")
\end{verbatim}

Possible output:

\begin{verbatim}
z:      42j
real:   0.0
imag:   42.0
type:   <class 'complex'>
\end{verbatim}

A complex number with no imaginary contribution can still be complex:

\begin{verbatim}
z = complex(42, 0)

print(f"z:      {z}")
print(f"real:   {z.real}")
print(f"imag:   {z.imag}")
print(f"type:   {type(z)}")
\end{verbatim}

Possible output:

\begin{verbatim}
z:      (42+0j)
real:   42.0
imag:   0.0
type:   <class 'complex'>
\end{verbatim}

The constructor \texttt{complex()} can build a complex object from two numeric parts:

\begin{verbatim}
real_part = 3
imag_part = 4

z = complex(real_part, imag_part)

print(f"z: {z}")
print(f"type(z): {type(z)}")
print(f"z.real: {z.real}")
print(f"z.imag: {z.imag}")
\end{verbatim}

Possible output:

\begin{verbatim}
z: (3+4j)
type(z): <class 'complex'>
z.real: 3.0
z.imag: 4.0
\end{verbatim}

The method \texttt{conjugate()} returns a related complex number where the sign of the imaginary part is reversed:

\begin{verbatim}
z = 3 + 4j
z_conj = z.conjugate()

print(f"z:           {z}")
print(f"conjugate:   {z_conj}")

print()
print(f"z real:      {z.real}")
print(f"z imag:      {z.imag}")
print(f"conj real:   {z_conj.real}")
print(f"conj imag:   {z_conj.imag}")
\end{verbatim}

Possible output:

\begin{verbatim}
z:           (3+4j)
conjugate:   (3-4j)

z real:      3.0
z imag:      4.0
conj real:   3.0
conj imag:   -4.0
\end{verbatim}

This method does not mutate the original complex object. Complex numbers are immutable. It returns a new complex object.

\begin{verbatim}
z = 3 + 4j

old_id = id(z)
z_conj = z.conjugate()

print(f"z:       {z}")
print(f"z_conj:  {z_conj}")

print()
print(f"id(z):      {id(z)}")
print(f"id(z_conj): {id(z_conj)}")
print(f"same object: {id(z) == id(z_conj)}")
\end{verbatim}

Possible output:

\begin{verbatim}
z:       (3+4j)
z_conj:  (3-4j)

id(z):      2214135700944
id(z_conj): 2214135700976
same object: False
\end{verbatim}

\subsubsection{Complex Arithmetic Boundaries: Why Ordering Comparisons Are Not Defined for Complex Numbers}

Complex numbers support equality comparison:

\begin{verbatim}
a = 3 + 4j
b = 3 + 4j
c = 3 - 4j

print(f"a == b: {a == b}")
print(f"a == c: {a == c}")
print(f"a != c: {a != c}")
\end{verbatim}

Possible output:

\begin{verbatim}
a == b: True
a == c: False
a != c: True
\end{verbatim}

However, Python does not define ordinary ordering for complex numbers:

\begin{verbatim}
a = 3 + 4j
b = 5 + 0j

print(a < b)
\end{verbatim}

This raises a \texttt{TypeError}.

Safe demonstration:

\begin{verbatim}
a = 3 + 4j
b = 5 + 0j

try:
    print(a < b)
except TypeError as err:
    print("ordering failed")
    print(f"message: {err}")
    print(f"type(err): {type(err)}")
\end{verbatim}

Possible output:

\begin{verbatim}
ordering failed
message: '<' not supported between instances of 'complex' and 'complex'
type(err): <class 'TypeError'>
\end{verbatim}

The reason is conceptual. Real numbers have a natural left-to-right order on a single number line. Complex numbers have two components. There is no single built-in rule that says whether \texttt{3 + 4j} should be less than \texttt{5 + 0j}.

A programmer can choose a comparison rule for a specific problem, but Python does not pretend that one universal rule exists.

For example, one possible rule is to compare by magnitude. Python's \texttt{abs()} returns the distance of the complex number from zero:

\begin{verbatim}
a = 3 + 4j
b = 5 + 0j
c = 1 + 1j

print(f"abs(a): {abs(a)}")
print(f"abs(b): {abs(b)}")
print(f"abs(c): {abs(c)}")

print()
print(f"abs(a) == abs(b): {abs(a) == abs(b)}")
print(f"abs(c) < abs(a):  {abs(c) < abs(a)}")
\end{verbatim}

Possible output:

\begin{verbatim}
abs(a): 5.0
abs(b): 5.0
abs(c): 1.4142135623730951

abs(a) == abs(b): True
abs(c) < abs(a):  True
\end{verbatim}

This comparison is not complex ordering. It is float ordering after reducing each complex number to a magnitude.

A list of complex values cannot be sorted directly:

\begin{verbatim}
values = [3 + 4j, 1 + 1j, 5 + 0j]

try:
    print(sorted(values))
except TypeError as err:
    print("sorting failed")
    print(f"message: {err}")
\end{verbatim}

Possible output:

\begin{verbatim}
sorting failed
message: '<' not supported between instances of 'complex' and 'complex'
\end{verbatim}

But it can be sorted by a chosen key, such as magnitude:

\begin{verbatim}
values = [3 + 4j, 1 + 1j, 5 + 0j]

sorted_by_magnitude = sorted(values, key=abs)

print(f"original:            {values}")
print(f"sorted by magnitude: {sorted_by_magnitude}")
\end{verbatim}

Possible output:

\begin{verbatim}
original:            [(3+4j), (1+1j), (5+0j)]
sorted by magnitude: [(1+1j), (3+4j), (5+0j)]
\end{verbatim}

The practical summary is:

\begin{quote}
Python's \texttt{complex} type is a built-in numeric object with floating-point real and imaginary parts. It supports complex arithmetic and equality checks, but it does not support ordinary ordering comparisons such as \texttt{<} or \texttt{>}. If ordering is needed, the programmer must explicitly choose what property is being ordered.
\end{quote}